\chapter{Introducción y Arquitectura General}
\label{sec:introduccion}

La placa \texttt{ctrl\_pap} acciona un motor paso a paso bipolar de dos fases dentro de la máquina CNC de tres
ejes. Cada eje lleva un módulo idéntico que recibe órdenes simples del controlador principal, genera la secuencia
de pasos por hardware y conmuta la potencia sobre los bobinados.

\section{Propósito dentro del Sistema CNC}
\label{sec:proposito_cnc}

En el sistema general optamos por separar el controlador embebido de las etapas de potencia mediante placas
dedicadas por eje. Esto aísla fallas eléctricas entre canales, ordena el cableado y reduce la cantidad de pines
necesarios en el microcontrolador a solo tres salidas por eje: paso, sentido y habilitación.

Además, la placa saca hacia una bornera externa el retorno común de corriente de los dos puentes H. Ese terminal
conecta con el módulo de sensado analógico del CNC para medir la carga mecánica del eje y detectar topes o
sobrecargas.

\section{Diagrama en Bloques e Interfaces Externas}
\label{sec:bloques_interfaces}

El circuito se organiza en cinco bloques: protección y regulación de tensión, adaptación de entrada, secuenciador
digital, puente H con diodos de protección y retorno de corriente. En la figura~\ref{fig.ej1:bloques_placa} se
observa cómo se vinculan entre sí y con las cuatro borneras perimetrales de paso $5\,\unit{\milli\meter}$.

\begin{figure}[!ht]
  \centering
  \resizebox{\textwidth}{!}{\begin{tikzpicture}[
  >=stealth, thick,
  font=\footnotesize,
  box/.style={
    draw=blue!70!black, fill=blue!8, rectangle, rounded corners=4pt,
    minimum height=1.25cm, minimum width=2.7cm, align=center
  },
  pwr/.style={
    draw=orange!80!black, fill=orange!10, rectangle, rounded corners=4pt,
    minimum height=1.25cm, minimum width=3.1cm, align=center
  },
  term/.style={
    draw=green!60!black, fill=green!10, rectangle, rounded corners=3pt,
    minimum height=1.0cm, minimum width=2.0cm, align=center
  },
  sens/.style={
    draw=purple!70!black, fill=purple!10, rectangle, rounded corners=4pt,
    minimum height=1.1cm, minimum width=3.0cm, align=center
  }
]

  % Fila superior: Alimentacion
  \node[term] (j1) {Bornera J1\\{\scriptsize V+, GND}};
  \node[pwr, right=0.9cm of j1] (reg) {Protección y\\Regulación\\{\scriptsize D1, C1, C5, LM7805CT}};

  % Fila central: Cadena de senal
  \node[term, below=1.6cm of j1] (j3) {Bornera J3\\{\scriptsize STEP, DIR3V3, EN}};
  \node[box, right=0.8cm of j3] (ifz) {Interfaz y\\Nivelación\\{\scriptsize Q1, Q2, LED1}};
  \node[box, right=1.1cm of ifz] (seq) {Secuenciador\\Digital\\{\scriptsize 74LS175, CD4052}};
  \node[pwr, right=1.3cm of seq] (drv) {Etapa de Potencia\\y Protección\\{\scriptsize L298N, D2--D9}};
  \node[term, right=0.8cm of drv] (j4) {Bornera J4\\{\scriptsize OUT1--OUT4}};

  % Fila inferior: Retorno de corriente
  \node[sens, below=1.3cm of drv] (ret) {Retorno Común\\{\scriptsize SENSE\_A (Pines 1 y 15)}};
  \node[term, right=0.8cm of ret] (j5) {Bornera J5\\{\scriptsize SENSE\_A, GND}};

  % Contenedor de placa
  \draw[dashed, draw=gray!70, rounded corners=6pt]
    ([xshift=-0.25cm, yshift=0.35cm]reg.north west -| ifz.west)
    rectangle
    ([xshift=0.25cm, yshift=-0.35cm]ret.south east -| drv.east);
  \node[anchor=north east, font=\scriptsize\bfseries, text=black!75]
    at ([xshift=0.15cm, yshift=0.25cm]reg.north east -| drv.east) {Placa ctrl\_pap};

  % Conexiones
  \draw[->] (j1) -- (reg);
  \draw[->, draw=orange!80!black] (reg.east) -|
    node[pos=0.25, above, font=\scriptsize] {VS} (drv.north);
  \draw[->, draw=blue!70!black] (reg.south) |- ([yshift=0.5cm]seq.north)
    node[pos=0.72, above, font=\scriptsize] {VSS ($+5\V$)} -- (seq.north);
  \draw[->, draw=blue!70!black] ([yshift=0.5cm]seq.north) -| (ifz.north);

  \draw[->] (j3) -- (ifz);
  \draw[->] (ifz) -- node[above, font=\scriptsize] {DIR} (seq);
  \draw[->] (ifz.south) -- ++(0,-0.45) -|
    node[pos=0.25, below, font=\scriptsize] {STEP} ([xshift=-0.3cm]seq.south);
  \draw[->] ([yshift=-0.45cm]ifz.south) -- ++(0,-0.35) -|
    node[pos=0.32, below, font=\scriptsize] {ENABLE} ([xshift=-0.5cm]drv.south);
  \draw[->] (seq) -- node[above, font=\scriptsize] {IN1--IN4} (drv);
  \draw[->] (drv) -- (j4);
  \draw[->] (drv) -- (ret);
  \draw[->] (ret) -- (j5);
\end{tikzpicture}
}
  \caption{Diagrama en bloques funcional de la placa \texttt{ctrl\_pap} y sus borneras de interfaz.}
  \label{fig.ej1:bloques_placa}
\end{figure}

En la tabla~\ref{tab.ej1:conectores} se detalla el pinado completo de las borneras \texttt{J1}, \texttt{J3},
\texttt{J4} y \texttt{J5}.

\begin{table}[!ht]
  \centering
  \resizebox{\textwidth}{!}{
    \begin{tabular}{|c|c|l|p{9.0cm}|}
      \hline
      \textbf{Ref.} & \textbf{Pin} & \textbf{Red} & \textbf{Función en el circuito} \\ \hline
      \multirow{2}{*}{\texttt{J1}} & 1 & \texttt{V+} &
        Entrada positiva de alimentación de potencia ($12\V$ a $24\V$ típicos). \\ \cline{2-4}
       & 2 & \texttt{GND} & Referencia común de masa de alimentación. \\ \hline
      \multirow{4}{*}{\texttt{J3}} & 1 & \texttt{GND} &
        Masa de referencia para las señales lógicas de control. \\ \cline{2-4}
       & 2 & \texttt{ENABLE} &
        Habilitación simultánea de ambos puentes H en \texttt{IC1} (activa en alto). \\ \cline{2-4}
       & 3 & \texttt{DIR3V3} & Selección del sentido de giro proveniente de lógica de $3.3\V$. \\ \cline{2-4}
       & 4 & \texttt{STEP} & Tren de pulsos de paso; avanza un paso en cada flanco ascendente. \\ \hline
      \multirow{4}{*}{\texttt{J4}} & 1 & \texttt{OUT1} &
        Terminal positivo del devanado de fase A del motor. \\ \cline{2-4}
       & 2 & \texttt{OUT2} & Terminal negativo del devanado de fase A del motor. \\ \cline{2-4}
       & 3 & \texttt{OUT3} & Terminal positivo del devanado de fase B del motor. \\ \cline{2-4}
       & 4 & \texttt{OUT4} & Terminal negativo del devanado de fase B del motor. \\ \hline
      \multirow{2}{*}{\texttt{J5}} & 1 & \texttt{GND} &
        Masa de referencia del retorno de corriente. \\ \cline{2-4}
       & 2 & \texttt{SENSE\_A} &
        Retorno unificado de los emisores inferiores de \texttt{IC1} hacia sensado externo. \\ \hline
    \end{tabular}
  }
  \caption{Asignación de pines en las borneras externas de la placa \texttt{ctrl\_pap}.}
  \label{tab.ej1:conectores}
\end{table}

El funcionamiento sigue una cadena directa: al alimentar \texttt{J1}, la placa alimenta los puentes de potencia y
regula $+5\V$ para su propia lógica. Luego, por \texttt{J3} entran las señales de $3.3\V$: \texttt{ENABLE}
habilita las salidas, \texttt{DIR3V3} se eleva a $5\V$ para fijar el sentido y cada pulso en \texttt{STEP}
actualiza las cuatro líneas \texttt{IN1}--\texttt{IN4} que invierten las bobinas en \texttt{J4}.
